\section{Evidencia del *chunking* dinámico: cuándo los límites de aprendizaje superan al *token* codificado}

El mecanismo de H-Net parece razonable, pero "aprendible" no equivale automáticamente a "mejor". Este capítulo distingue tres conjuntos de evidencia para diferentes problemas: el primero pregunta si un *learned chunk* supera al BPE tras suficiente datos; el segundo indaga por qué el *outer stage* prefiere Mamba incluso cuando la entrada ya es BPE; y el tercero sitúa los métodos en el ADN, observando si la pendiente de *scaling* cambia sin una *heuristic* del *tokenizer* madura. Al leer la figura se deben considerar simultáneamente el cómputo emparejado (*matched compute*), el punto de cruce (*crossover*) y el límite del dominio.

\subsection{Punto de cruce con cómputo emparejado: "ganar después de perder" es la forma típica de la lección amarga}

En el eje horizontal de la figura siguiente se muestran los bytes totales de entrenamiento; en el vertical, los bits por byte. El modelo FLOPs y la escala aproximada están ajustados. Los Transformers negros utilizan BPE directamente; una capa de H-Net con *byte* es peor al inicio porque debe aprender simultáneamente el lenguaje y los límites; tras suficientes datos, las curvas se cruzan. Dos capas de H-Net aprenden aún más abstracciones multinivel, pero introducen más parámetros y dificultad de optimización.

\lecturefigure{slide-027.jpg}{*Scaling* del *chunking* dinámico: comparación con cómputo emparejado de BPE isobitrópico, H-Net de una sola capa y de múltiples capas}{Video oficial de Stanford, 00:40:40--00:44:09}

\begin{knowledgebox}{Leer la figura: tres pasos para verificar el *crossover*}
Primero, confirmar si las curvas se comparan bajo presupuestos de FLOP aproximadamente iguales, en lugar de solo coincidir parámetros; segundo, observar la desventaja del modelo con *learned chunk* en la zona de pocos datos, reconociendo que el BPE diseñado manualmente proporciona un sesgo inductivo fuerte; tercero, fijarse en la pendiente y el punto de cruce en la zona de muchos datos para juzgar si el *scaling* adicional compensa el costo de aprender los límites.
\end{knowledgebox}

Este tipo de experiencia de *scaling* puede escribirse como:
\[
\mathcal{L}(C)\approx \mathcal{L}_\infty + aC^{-\alpha},
\]
donde $C$ es el cómputo de entrenamiento o un presupuesto de datos aproximadamente equivalente, $\mathcal{L}_\infty$ es la pérdida irreductible, $a$ es una constante y $\alpha$ es el exponente de *scaling*. Cuando se dice en la figura que "escala mejor", no se refiere a un punto más bajo, sino a que, a medida que aumenta $C$, la curva de la jerarquía aprendida mantiene una tendencia descendente más favorable. El ajuste en intervalos finitos aún puede verse afectado por el optimizador, la mezcla de datos y el dimensionamiento del modelo.

\teachervoice{00:42:37--00:44:21, Gu explica el *crossover* con la lección amarga: las características manuales ayudan mucho en un régimen limitado por datos; el *learned chunk* necesita gastar datos primero para aprender a segmentar, pero una vez que aprende unidades más adecuadas a la tarea que el BPE, puede seguir beneficiándose del *scaling*.}

\begin{warningbox}{Aprender "tokens" no garantiza extrapolación infinita}
Las curvas actuales solo cubren intervalos limitados de escala y datos. Jerarquías más profundas pueden saturarse prematuramente por colapso de enrutamiento, asignación de parámetros, inestabilidad de optimización o eficiencia del sistema. Para juzgar si el siguiente orden de magnitud seguirá siendo dominante, es necesario escalar verdaderamente los modelos y los datos, no prolongar la línea recta fuera del gráfico.
\end{warningbox}

\subsection{Ablación más crítica: incluso cuando la entrada ya es BPE, el *outer Mamba* sigue siendo importante}

Si la ventaja de los SSM proviene solo de que "la secuencia de bytes es demasiado larga", entonces tras cambiar la entrada a tokens BPE, debería ser suficiente usar un Transformer como codificador externo. Sin embargo, la figura muestra que: incluso cuando todo el sistema opera con BPE, añadir Mamba en las etapas externas mejora notablemente los BPB. Aquí, la tarea del *outer stage* no es simplemente leer la entrada, sino crear representaciones comprimibles para el enrutamiento; por ello, el sesgo inductivo de estado finito puede ayudar directamente a la formación de abstracciones.

\lecturefigure{slide-028.jpg}{Ablación del *outer stage* de H-Net con Mamba sigue mostrando un sesgo inductivo comprimido sobre entrada BPE}{Video oficial de Stanford, 00:44:10--00:46:09}

\begin{importantbox}{La evidencia más fuerte de esta clase de "la compresión es una característica"}
Cuando la unidad de entrada ya es BPE, la explicación por resolución no basta para justificar el beneficio del *outer Mamba*. La hipótesis más razonable es que el codificador previo al enrutamiento necesita integrar la historia local en representaciones aptas para el *chunking*, y el estado recurrente finito fuerza a la red a formar un resumen. Este resultado sigue siendo una ablación específica de H-Net, pero apoya más directamente "la compresión como sesgo inductivo" que los simples benchmarks de bytes.
\end{importantbox}

\begin{warningbox}{La ablación no prueba por sí sola el mecanismo causal único}
El *outer stage* con Mamba y Transformer pueden diferir en parametrización, normalización, optimización, kernel o profundidad efectiva. Los resultados indican que "los sistemas que usan Mamba son mejores", pero no prueban por sí solos que todo el beneficio provenga de la compresión de estado finito. Una verificación más estricta del mecanismo requiere arquitectura emparejada y *probe* de representación.
\end{warningbox}

\subsection{dnaHNet: cuando no hay un *tokenizer* semántico, la jerarquía cambia la pendiente de *scaling*}

El ADN no cuenta con una *heuristic* de segmentación madura como los espacios en blanco, raíces y frecuencias de palabras del inglés. Cada punto en la figura corresponde a un presupuesto diferente de FLOP y, bajo ese presupuesto, se varía (*sweep*) la asignación modelo/datos. La curva de H-Net no solo se traslada verticalmente, sino que muestra pendientes diferentes; el orador concluye que la jerarquía aprendida encuentra patrones que los *tokenizers* comerciales no pueden ofrecer.

\lecturefigure{slide-029.jpg}{Leyes de *scaling* dnaHNet: tendencia de crecimiento del *chunking* dinámico sin una *heuristic* madura de tokenización}{Video oficial de Stanford, 00:47:20--00:49:19}

Esta clase utiliza la versión v3 del artículo arXiv de `dnaHNet` (2026-04-09) según la fecha del curso, ya que es la última versión disponible antes del 2026-04-16. La figura cita el preentrenamiento autoregresivo genómico; no se debe traducir directamente una pérdida menor a comprensión causal biológica. El valor real en tareas downstream también depende de la validación en elementos reguladores, efectos variantes y contexto de tipo celular.

\begin{knowledgebox}{Por qué el ADN es una prueba más fuerte que el inglés basado en caracteres}
Aunque el inglés permite modelar bytes, ya cuenta con un BPE maduro; los investigadores siempre pueden preguntar "¿por qué no usar un *tokenizer*?". La secuencia de bases unitarias del ADN carece de límites universales de palabras similares; definir tokens diseñados manualmente es más difícil. Si el *learned chunking* cambia el *scaling* aquí, se acerca más al problema original de datos que pretendía resolver.
\end{knowledgebox}

\teachervoice{00:47:17--00:49:18, el orador llama al inglés basado en caracteres un a largo plazo con carácter académico, mientras considera el ADN como el escenario donde la tokenización realmente no puede ser semantizada; esto también es la razón por la que cree que el *chunking* dinámico tiene más relevancia práctica.}

\subsection{Resumen de la sección}

Las evidencias de H-Net forman una cadena progresiva: los *chunks* aprendidos a nivel de byte superan al BPE tras suficientes datos; las ablaciones en etapas externas sobre entrada BPE aún prefieren Mamba, apuntando a un sesgo inductivo por compresión; y el ADN muestra diferencias más fuertes en *scaling* cuando falta un *tokenizer* semántico. La cadena apoya "la compresión puede ayudar a formar abstracciones", pero sigue limitada por escala finita, coincidencia de arquitectura y transferencia de dominio.
